\documentclass[11pt]{article}

% ---- Overleaf-ready preamble (compiles with pdfLaTeX; no shell-escape needed) ----
\usepackage[margin=1in]{geometry}
\usepackage{booktabs}
\usepackage{amsmath}
\usepackage{xcolor}
\usepackage{listings}
\usepackage[hidelinks]{hyperref}

\definecolor{codebg}{gray}{0.96}
\lstdefinestyle{json}{
  basicstyle=\ttfamily\small,
  backgroundcolor=\color{codebg},
  breaklines=true,
  frame=single,
  framesep=4pt,
  showstringspaces=false,
  columns=fullflexible,
  xleftmargin=2pt,
}

\title{PennyPilot: Permission-Gated, Hidden-Constraint Shopping with\\
Reinforcement-Learning Post-Training on Constrained Hardware}
\author{Param Madan}
\date{\today}

\begin{document}
\maketitle

\begin{abstract}
We present PennyPilot, a multi-turn, multilingual shopping agent trained via a
full reinforcement-learning post-training pipeline (SFT $\rightarrow$ RLOO/GRPO)
on a single 32\,GB V100 GPU. The task is designed so that naive strategies fail:
the shopper's budget and 2--3 hard requirements are hidden and revealed only
through correct clarifying questions, and under partial discovery the cheapest
visible item is a distractor in 99/100 scenarios ($n{=}100$, measured). Rewards
are computed against catalog ground truth --- there is no learned reward model,
and therefore nothing to hack. A structural permission gate makes any cart
addition without an explicit, per-SKU grant unrewardable, holding safety at 0
violations across every evaluation run. We report the full progression from a
1.5B validation model to a 7B production model, a pre-registered comparison of
RLOO against GRPO under matched rollout budgets, a systems/efficiency campaign
(SS0--SS13) that measured rather than assumed every optimization, and ---
critically --- the negative and partial results alongside the positive ones.
PennyPilot is the post-training arm of a longer lineage: it extends
\emph{shoprl-fabric}, a single-machine multi-turn RL platform, by adding hidden
needs, a permission gate, and an information-gain reward.
\end{abstract}

\section{Introduction}
Most public RL-post-training demonstrations either (a) train on tasks that are
already near-saturated for the base model, making the ``improvement'' hard to
distinguish from noise, or (b) omit the systems engineering that determines
whether the loop is even trainable on the hardware actually available to most
practitioners (a single consumer/academic GPU, not a multi-node cluster).
PennyPilot addresses both: the task's difficulty is \emph{measured and gated}
before any training spend, and the entire pipeline --- rollout, reward,
optimization, evaluation, checkpointing --- is built and validated on a single
Volta-generation V100.

The agent is the workload; the actual engineering target is the training
infrastructure: multi-turn RL stability, a safety mechanism that survives
optimization pressure, and profiling-driven rollout efficiency.

\section{Task Design}
The shopper's true objective is hidden from the agent: a budget ceiling and 2--3
hard constraints (brand, minimum RAM, minimum battery life, maximum weight),
disclosed one at a time and only in response to a relevant clarifying question.
This is a deliberate difficulty mechanism, not decoration: under partial
constraint discovery, the cheapest catalog item satisfying the \emph{known}
constraints fails at least one \emph{undiscovered} constraint in 99 of 100
generated scenarios. A policy that grabs the first plausible result loses; a
policy that disciplines itself to fully enumerate constraints before committing
wins. This gives the reward signal real texture instead of a near-ceiling pass
rate from turn one.

Rewards are scored against the catalog's real specifications --- never against
the model's stated beliefs about them --- so there is no reward model in the loop
and no surface for reward hacking via plausible-sounding but false claims.

\section{Method}

\subsection{Structured action interface}
The policy emits single-line JSON actions rather than raw browser selectors or
free-form tool calls:
\begin{lstlisting}[style=json]
{"action": "ask_user", "question": "What is your total budget?"}
{"action": "search", "query": "laptop under 2000 dollars"}
{"action": "select_product", "product_id": "LAP-0014", "reason": "cheapest valid"}
{"action": "request_cart_permission", "items": ["LAP-0014"], "estimated_total": 796.29}
{"action": "add_to_cart", "product_id": "LAP-0014"}
\end{lstlisting}
Environment adapters translate these into effects: a synthetic catalog
environment (training and most evaluation), a WebShop-compatible adapter
(unseen-product generalization test), and a Playwright browser mirror used only
for human-facing demo/replay --- the browser is never in the training loop, only
a projection of decisions already made.

\subsection{Structural permission gate}
\texttt{add\_to\_cart} for a SKU succeeds only if that exact SKU appears in an
explicit, previously granted permission set. This is enforced in the
environment's dialogue-state transition function, \emph{not} in the reward term:
no combination of policy weights can produce a rewarded cart addition without a
real, matching grant having occurred first. Ambiguous replies do not count as
approval, and a held ``no'' survives subsequent permission requests for the same
or a different item until explicitly lifted. Measured result: 0 permission
violations across every SFT, RLOO, GRPO, and zero-shot-transfer evaluation run in
this work.

\subsection{Environment-side dialogue state}
Corrections (``actually, my budget is \$90,'' ``I already have sunscreen'') are
applied to a structured dialogue state that invalidates stale plans; the model's
context window is a \emph{view} onto this state, never its source of truth. A
deterministic bilingual (English/Spanish, including code-switched input) language
layer extracts constraints regardless of which language they arrive in.

\subsection{Loss-mask verification}
SFT loss is computed on policy tokens only. A reference-exact verifier crashes the
training run if any environment- or user-authored token ever leaks into the loss
--- the specific failure mode in which training loss improves while the policy's
actual behavior degrades, because it is partially learning to imitate the
environment rather than to act in it.

\subsection{Precision on Volta}
The V100 (Volta) lacks native \texttt{bf16} tensor-core support; naive
\texttt{bf16} training on this hardware falls back to software emulation. We
instead train in \texttt{fp16} with dynamic loss scaling (\texttt{GradScaler}),
which is necessary, not optional --- unscaled \texttt{fp16} training diverges at
step 1 on this task. Measured: \texttt{fp16}+\texttt{GradScaler} is
$4.27\times$ faster than the \texttt{bf16}-emulated alternative, with both
numerically stable once the scaler is applied.

\section{Experimental Setup}
\textbf{Models.} Qwen2.5-1.5B-Instruct is used first, to validate the entire
pipeline (environment, SFT recipe, RL loop, evaluation harness) at low cost.
Qwen2.5-7B-Instruct is the production line. For RL, a shared-reference trick ---
one frozen base model carrying two separate LoRA adapters (policy and reference)
--- lets 7B policy+reference RL fit in 32\,GB, avoiding two full model copies.

\textbf{Hardness gate.} Before any training spend, we measure the untrained base
model's success rate on the held-out hard split ($n{=}64$): 12.5\%. This falls
inside a pre-registered acceptance band of 10--40\% --- enough signal to
bootstrap SFT from, enough headroom for RL to matter. A task solved at 90\%
zero-shot would not have justified this pipeline.

\textbf{Algorithms.} RLOO (critic-free, leave-one-out baseline) is the bring-up
algorithm; GRPO is compared against it at an \emph{equal} rollout budget (400
trajectories) so that any difference is attributable to the advantage rule, not
to sample count.

\section{Results}
\begin{table}[h]
\centering
\begin{tabular}{lc}
\toprule
\textbf{Stage} & \textbf{Task success (es-en hard split)} \\
\midrule
Base model (hardness gate) & 12.5\% \\
SFT (env-recorded multilingual demos) & 48.4\% (RL-difficulty); 98--100\% (standard) \\
RLOO, 50 steps & 92.2\% \\
GRPO, equal 400-trajectory budget & 100.0\% \\
GRPO + 4-epoch clipped reuse & 100.0\% (KL max 0.027) \\
\bottomrule
\end{tabular}
\caption{Learning progression. Every arm is violation-free and every number is
held-out.}
\end{table}

\textbf{H1 --- GRPO instability is regime-dependent, not intrinsic.} GRPO's
reputation for training instability is often reported without reference to the
reward regime it was observed in. Here, measured KL stays at 0.013 (a
healthy-variance regime) versus 0.58--0.99 drift on a saturated task using the
identical stack. Both the hypothesis and the acceptance threshold were registered
before the run.

\textbf{H2 --- multi-epoch reuse is sufficient but not necessary.} The reuse arm
(H2) reached the same 100\% ceiling as single-pass advantage scaling; reuse did
not add task success beyond what scaling alone already achieved, and was
effectively free in wall-clock terms since rollout dominates iteration time
regardless (Sec.~\ref{sec:systems}). This mechanism-level finding, not just the
top-line number, was the pre-registered target.

\textbf{H3 --- catastrophic forgetting is rehearsal-mitigable at no cost (7B).} A
specialist SFT arm (task-only demonstrations) loses general-chat retention
(12.5\%$\rightarrow$10.0\%) alongside modest shopping competence (0.820). A
rehearsal arm --- 25\% self-distilled general-chat data mixed into the same SFT
stage --- holds retention at 97.5\% while \emph{exceeding} the specialist's
shopping competence (0.859) after identical RL. Rehearsal was not a trade-off
here; it dominated on both axes.

\textbf{Zero-shot generalization (WebShop).} Transferred to an unseen-product,
different-backend environment with no fine-tuning: 0 permission violations
(safety transfers completely), 60\% task success (competence transfers
partially). This is reported as an honest domain-shift datapoint, not elided.

\textbf{Preemption safety.} A live mid-run SIGTERM $\rightarrow$ Slurm requeue
$\rightarrow$ resume drill produced a clean 20-of-20-step continuation with KL
trace continuous across the kill seam --- preemption safety was measured
operationally, not only asserted from the checkpoint code.

\section{Systems and Efficiency}
\label{sec:systems}
A 14-row profiling campaign (SS0--SS13) measured every proposed optimization
before adopting or rejecting it:
\begin{itemize}
\item Rollout generation is 95.8\% of RL iteration wall-clock (three consistent
profiling passes) --- the correct place to spend effort is rollout, not the
optimizer step.
\item End-to-end iteration time: 93.2\,s $\rightarrow$ 13.4\,s ($6.96\times$),
via batched lockstep rollouts ($4.98\times$ more tokens/iteration, $6.4\times$
more episodes/minute), merge-at-sync vLLM serving ($2.3\times$ alone --- and the
workaround for Volta's non-functional LoRA-serving kernels), and checkpoint-off
updates at small batch sizes ($1.48\times$).
\item Measured and \emph{rejected}: an inference-engine swap alone ($1.06\times$,
not worth its complexity in isolation), prefix caching (hardware-gated on Volta
and $\le 7\%$/turn at these sequence lengths even where available), request
packing (8.2\% token waste at this distribution), and a custom Triton kernel ---
profiled ceiling of 0.84\% of iteration time, and therefore deliberately not
written. The profiling data is the justification for \emph{not} building
something, which is as load-bearing a result as the optimizations that were
adopted.
\end{itemize}

\section{Limitations (honest)}
\begin{itemize}
\item The scripted user simulator has cosmetic phrasing quirks; a frozen-LLM user
model is a planned drop-in behind the same interface.
\item Mid-flow ground-truth corrections (e.g., a live budget revision) are
deferred pending scenario-schema support --- not implemented, not faked.
\item The WebShop evaluation uses a dialect-faithful in-repo backend, verified
against the real WebShop parser; the original princeton-nlp instance remains a
documented, not-yet-wired backend.
\item \texttt{vLLM} LoRA serving and prefix caching are hardware-gated on Volta
(the relevant Triton kernels require Ampere or newer); the merge-at-sync
workaround is documented with its correctness asserts.
\item H2's registered mechanism (multi-epoch reuse) was not the reason GRPO
reached its ceiling --- advantage scaling alone was sufficient. The reuse arm's
\emph{own} hypothesized benefit is therefore an open, documented follow-up.
\item Reward-shaping extension terms specified in early design notes
(need-coverage, state-consistency, trajectory-efficiency) were not implemented;
the shipped reward is the verifiable base term plus information gain. Its visible
cost is a measured tendency toward over-asking clarifying questions --- an honest,
named consequence rather than an unexplained quirk.
\item No preference-optimization stage (DPO/RLHF-style) has been run, despite
collecting explicit thumbs-up/down feedback during live use that would support
one. This is the single largest remaining gap between this work and a
``complete'' modern post-training pipeline.
\end{itemize}

\section{Vision: Agents over a Self-Hosted Distributed Retrieval Substrate}
PennyPilot's \texttt{search} action is served, today, by a synthetic catalog
environment --- a controlled ground-truth index that makes rewards verifiable.
The natural next step is to replace that index with a real retrieval substrate
without giving up the properties that made the training loop honest.

We are building that substrate separately as \emph{Aether}: a from-scratch,
single-operator \textbf{distributed spatial--vector search engine in Rust}
(sharded \texttt{shard-node}s behind a \texttt{coordinator}, gRPC transport,
leader election with automatic failover, cross-cluster federation, a
natural-language query layer, and provenance-carrying briefings). Aether is
already operational end to end --- a running cluster demonstrates coordinator-
driven leader failover and query fan-out across shards.

The vision is their composition: PennyPilot as the \emph{policy} and Aether as
the \emph{environment/retrieval tier}. The agent's structured \texttt{search}
action fans out over Aether's sharded index instead of a synthetic catalog;
constraint filters become Aether's structured filters (which already aggregate
without shipping raw values); and the agent's final briefing inherits Aether's
provenance trail, so a recommendation is traceable to the exact shards and
records that produced it. This turns a controlled RL benchmark into a
self-hosted agentic-retrieval system --- an LLM policy that reasons and asks,
grounded on a fault-tolerant distributed index that a single practitioner can
own and operate. Making the reward verifiable against a \emph{live, distributed}
ground truth (rather than a static synthetic one) is the central open research
problem this composition raises, and the direction we intend to pursue.

\section{Conclusion}
PennyPilot demonstrates that rigorous RL post-training --- pre-registered
hypotheses, measured hardness gating, honest negative and partial results, and a
safety mechanism verified rather than assumed --- is achievable end to end on a
single academic GPU. The central methodological claim is not any individual
number in Section~5, but the discipline connecting them: every reported result
was predicted before it was measured, and every reported optimization was adopted
or rejected based on a profile, not intuition.

\end{document}
